\documentclass{elsarticle}
\usepackage{amsmath}
\usepackage{amssymb}
\usepackage{graphicx}
\usepackage{booktabs}
\usepackage{siunitx}
\usepackage{chemfig}
\usepackage{mhchem}
\usepackage{bpchem}
\usepackage{hyperref}
\usepackage{url}

\sisetup{
  detect-all,
  per-mode=symbol
}

\graphicspath{{./figures/}}

\begin{frontmatter}

\title{Synthesis and characterization of novel metal-organic frameworks for carbon capture}

\author[aff1]{First Author}
\address[aff1]{Department of Chemistry, Example University, City, Country}
\author[aff2]{Second Author}
\address[aff2]{Institute of Materials Science, Research Institute, Country}

\begin{abst}
\begin{abstract}
Metal-organic frameworks (MOFs) have emerged as promising materials for
carbon capture applications~\cite{li2012metal, furukawa2013water}.
We report the synthesis of a novel Zr-based MOF with exceptional
CO$_2$ adsorption capacity. The material demonstrates a record-high
CO$_2$ uptake of \SI{8.5}{\mmol\per\gram} at \SI{298}{\kelvin}
and \SI{1}{\bar}.
\end{abstract}
\end{abst}

\keywords{Metal-organic framework; Carbon capture; CO$_2$ adsorption; Zirconium; Post-combustion}

\end{frontmatter}

\section{Introduction}

Climate change driven by anthropogenic CO$_2$ emissions necessitates
the development of efficient carbon capture technologies~\cite{gibson}.

\section{Results and Discussion}

The crystal structure was determined by single-crystal X-ray diffraction.
The framework exhibits a cubic topology with a pore size of \SI{12.5}{\angstrom}.

\subsection{Synthesis}

The MOF was synthesized by solvothermal reaction:
\begin{equation}
\ce{ZrCl4 + H4L ->[DMF, 393 K] MOF}
\end{equation}

The ligand H$_4$L was prepared according to the literature procedure.

\subsection{Gas Adsorption}

CO$_2$ adsorption isotherms were measured at various temperatures:

\begin{table}[ht]
\centering
\caption{CO$_2$ adsorption capacity at different conditions.}
\label{tab:co2}
\begin{tabular}{lS[table-format=2.1]S}
  \toprule
  Temperature (\si{\kelvin}) & \multicolumn{2}{c}{Uptake (\si{\mmol\per\gram})} \\
  \cmidrule{2-3}
   & \SI{0.15}{\bar} & \SI{1}{\bar} \\
  \midrule
  298 & 2.3 & 8.5 \\
  323 & 1.8 & 6.9 \\
  373 & 1.2 & 4.8 \\
  \bottomrule
\end{tabular}
\end{table}

\section{Conclusion}

We have developed a novel Zr-based MOF with outstanding CO$_2$ capture
performance. The material shows great promise for post-combustion carbon capture.

\section*{Acknowledgments}

This work was supported by the National Science Foundation.

\bibliographystyle{elsarticle-num}
\bibliography{refs}

\end{document}
